\section{Beam-Conditioning Simulation Framework}

This framework models a laser beam propagated through apertures, optional lenses,
and optional spatial-filter elements. The purpose is to evaluate whether the beam
can be delivered into a \SI{3}{mm} diameter target region at a distance of
\SI{1}{m}. The primary metric is the fraction of transmitted optical power
outside a radius
\[
    r_{\rm target}=\SI{1.5}{mm}
\]
at
\[
    z_{\rm target}=\SI{1000}{mm}.
\]
The secondary metric is throughput. Configurations with high target containment
but very low throughput should be rejected.

\subsection{Files}

The code has been split into a reusable library and short demo scripts:
\begin{itemize}
    \item \texttt{beam\_conditioning\_core.py}: reusable propagation, lens,
    aperture, plotting, and optimization tools;
    \item \texttt{student\_common.py}: default input beam and shared helpers;
    \item \texttt{demo\_single\_propagation.py}: one fixed optical configuration;
    \item \texttt{demo\_subsystem\_comparison.py}: compares increasingly
    complicated systems;
    \item \texttt{demo\_fixed\_aperture\_scan.py}: scans only aperture radii
    while keeping lens positions fixed;
    \item \texttt{demo\_optimize.py}: advanced exploratory optimizer.
\end{itemize}

The recommended starting point is
\begin{verbatim}
python demo_single_propagation.py
\end{verbatim}
This produces target-plane plots and a multi-page PDF showing the beam at saved
planes. The PDF should be inspected to determine where the beam first becomes
broader, clipped, or structured.

\subsection{Model}

The code propagates a scalar complex field
\[
    U(x,y;z),
\]
with intensity
\[
    I(x,y;z)=|U(x,y;z)|^2.
\]
Propagation is performed using the angular-spectrum method. Lenses are treated
as thin-lens phase elements located at the appropriate principal planes computed
from the cemented-doublet first-order model. Apertures and pinholes are hard
circular stops. This model is useful for system-level diffraction and clipping
studies, but it is not a replacement for full CODE V or Zemax aberration
analysis.

\subsection{Recommended Study Plan}

The study should proceed in stages, with one clear question at each stage.

\paragraph{Stage 1: Single fixed propagation.}
Run \texttt{demo\_single\_propagation.py}. Record the throughput, fraction
outside the target radius, and RMS radius. Inspect the beam-plane PDF. The goal
is not to optimize, but to understand what each plane looks like.

\paragraph{Stage 2: Subsystem comparison.}
Run \texttt{demo\_subsystem\_comparison.py}. This compares:
\begin{enumerate}
    \item input beam only;
    \item input apertures only;
    \item input apertures plus lens 1;
    \item input apertures plus lens 1 and lens 2;
    \item lens pair with a pinhole;
    \item lens pair with a post-lens aperture.
\end{enumerate}
A configuration should not be considered useful unless it improves on the simpler
configuration before it.

\paragraph{Stage 3: Fixed-layout aperture scan.}
Run \texttt{demo\_fixed\_aperture\_scan.py}. This holds the lenses fixed in the
expected common-focus geometry and scans only aperture radii. This is deliberately
more constrained than a global optimization. It asks whether aperture-only cleanup
can reduce the light outside the \SI{3}{mm} target while keeping acceptable
throughput.

\paragraph{Stage 4: Input-beam robustness.}
After the baseline is understood, edit \texttt{student\_default\_input\_field}
in \texttt{student\_common.py}. Repeat the first three studies for a circular
Gaussian, elliptical Gaussian, tilted elliptical Gaussian, top-hat beam, and dirty
Gaussian. A useful configuration should not depend strongly on one idealized input
beam.

\paragraph{Stage 5: Numerical convergence.}
Promising cases should be rerun with larger grids. For example compare
\[
    N=1024,\quad L=\SI{18}{mm}
\]
with
\[
    N=2048,\quad L=\SI{18}{mm}.
\]
The target containment and encircled-energy curve should remain reasonably stable.

\paragraph{Stage 6: Advanced optimization.}
Only after the subsystem and aperture-only studies are understood should
\texttt{demo\_optimize.py} be used. The optimizer can find mathematically better
parameters, but these may not be physically meaningful. Any optimized candidate
must be validated by rerunning the subsystem comparisons and larger-grid checks.
